\begin{lemma}[Scheide maximum-triple fact (g)]
Every globally maximum triple $(T,e_0,\varphi)$ has at least four free
colors: $|\mathcal F(T,e_0,\varphi)|\geq4$.
\end{lemma}
\begin{proof}
Write $X=V(T)$.  By \noderef{ScheideGlobalMaximumOrder}, $X$ is
elementary.  The two
ends of $e_0$ are the first two distinct vertices of $T$ by
\noderef{TreeSequence}.  The exact counts in
\noderef{ScheideGlobalMaximumMissingCounts} say that each vertex of $X$
misses at least one color and that these two ends each miss one additional
color.  Elementariness makes the missing-color sets of distinct vertices
disjoint.  Therefore
\[
             |\overline\varphi(T)|\geq |X|+2.             \tag{1}
\]

If $T$ has $n$ tree edges, then it has $n+1=|X|$ vertices.  Only the
$n-1=|X|-2$ edge identities after the distinguished first edge can be
colored tree edges.  Consequently at most $|X|-2$ distinct colors are used
on $T$ in the sense of \noderef{TreeUsedColors}.  Removing all used colors
from the missing-color union in (1) leaves at least
$(|X|+2)-(|X|-2)=4$ colors.  This difference is exactly the free-color set
defined in \noderef{TreeFreeColors}.
\end{proof}
